\documentclass{uofa-eng-assignment}
\usepackage{caption}
\usepackage{hyperref}
\hypersetup{
    colorlinks = false,
    linkbordercolor = {white},
}
\captionsetup[table]{name=Tabla, labelsep=period}
\captionsetup[figure]{name=Figura, labelsep=period}

\newcommand{\hpl}{\hyperlink}
\newcommand{\hpt}{\hypertarget}
\newcommand*{\name}{Joan Sebastian Ramirez Sanchez  }
\newcommand*{\id}{202223948}
\newcommand*{\course}{Control óptimo aplicado a modelos biológicos}
\newcommand*{\assignment}{Tarea 3 de 3}
\newcommand*{\dated}{\today}
\newcommand{\tf}{\textbf}
\newcommand{\la}{\lambda}
\begin{document}
%Ejercicios propuestos 17.1, 17.3, 17.5, 20.2, 20.3
\maketitle
\section*{Ejercicio 17.1}
    Resolver
    \[
    \text{max} \{8x_1(18)+4x_2 (18)\}
    \]
    Sujeto a 
    \begin{align*}
        x_1'(t) =x_1(t)+x_2(t) , \quad x_1(0) = 15, \\
        x_2'(t) = 2 x_1(t)- u(t), \quad x_2(0)= 20, \\
        0 \leq u(t) \leq 1.
    \end{align*}
    Observe que la formula general del hamiltoneano se puede ver como 
    \[
    H = \lambda_0 L+ \lambda_1 f+ \cdots \lambda_nf_n
    \]
    En nuestro caso $L = 0$ por lo tanto tenemos que 
    \[
    H(t,x_1,x_2,\lambda_1,\lambda_2) = \lambda_1 (x_1 + x_2) + \lambda_2 (2x_1 - u)
    \]
    donde las ecuaciones adjuntas están dadas por 
    \begin{align*}
        \lambda'_1 &= - \frac{\partial H}{\partial x_1} = -(\lambda_1+2 \la_2) \quad \lambda'_2 = -\frac{\partial H}{\partial x_2} = -\la_1
    \end{align*}
    y
    \begin{align*}
        \psi (t) = \frac{\partial H}{\partial u} = -\lambda_2
    \end{align*}
    así si queremos maximizar el hamiltoneano $H$ con respecto a $u$ tenemos que hacer uso del signo de $\la_2$
    \begin{itemize}
        \item si $ \la_2 > 0$ entonces $-\la_2 <0$ entonces se elige $u$ mínimo $u*= 0$.
        \item si $\la_2 <0$ entonces $- \la_2 >0$ entonces se elige $u*=1$. 
        \item si $\la_2= 0$ entonces es un caso singular lo cual no ocurre
    \end{itemize}
    por lo tanto tenemos el siguiente sistema
    \begin{align*}
    \begin{pmatrix}
        \lambda_1' \\
        \lambda_2'
        \end{pmatrix}
        =
        \begin{pmatrix}
        -1 & -2 \\
        -1 & 0
        \end{pmatrix}
        \begin{pmatrix}
        \lambda_1 \\
        \lambda_2
        \end{pmatrix},
        \qquad
        \begin{pmatrix}
        \lambda_1(18) \\
        \lambda_2(18)
        \end{pmatrix}
        =
        \begin{pmatrix}
        8 \\
        4
        \end{pmatrix}.
    \end{align*}
    con valores propios $v_1 = 1, \quad v_2=-2$ y vectores propios $(1,-1)^T$ y $(2,1)^T$ respectivamente, la solución general:
    \[
    \la_1(t) = c_1 e^t+2c_2e^{-2t}, \quad \la_2 (t)= -c_1e^t+c_2e^{-2t}
    \] 
    aplicando las condiciones finales nos queda que 
    \[
    c_1 e^{18} + 2c_2 e^{-36} = 8,
    \qquad
    -c_1 e^{18} + c_2 e^{-36} = 4.
    \]
    Sumando ambas ecuaciones
    \[
    3c_2 e^{-36} = 12
    \;\Rightarrow\;
    c_2 = 4e^{36}.
    \]
    Sustituyendo: 
    \[
    -c_1 e^{18} + 4 = 4
    \;\Rightarrow\;
    c_1 = 0.
    \]
    Por lo tanto
    \[
    \lambda_1(t) = 8e^{36-2t},
    \qquad
    \lambda_2(t) = 4e^{36-2t}.
    \]
    Como $\la_2 >0$ para todo $t$, el control óptimo es $u*(t) = 0$, asi el sistema nos queda 
    \begin{align*}
        x'_1 = x_1+x_2, \quad x'_2 = 2x_1
    \end{align*}
    con las mismas condiciones y viéndolo de forma matricial  $x'=Ax$ con $A = \begin{pmatrix}
        1 & 1 \\
        2 & 0 
    \end{pmatrix}.
    $
    cuyos valores propios son $\{2,-1\}$ y vectores propios $\{(1,1)^T , (1,-2)^T\}$ respectivamente, por lo tanto la solución queda 
    \begin{align*}
        x_1(t) = d_1 e^{2t}+d_2 e^{-t}, \quad x_2(t) = d_1 e ^{2t} -2d_2 e^{-t} 
    \end{align*}
    con las condiciones iniciales 
    \[
    d_1+d_2= 15, d_1 -2d_2= 20 \Rightarrow d_1 = \frac{50}{3}, d_2= -\frac{5}{3}.
    \]
    Asi la solución es :
    \[
    x_1(t) = \frac{50}{3} e^{2t} -\frac{5}{3}e^{-t}, \quad x_2(t) = \frac{50}{3} e ^{2t} +\frac{10}{3} e^{-t} 
    \]
    y el funcional objetivo
    \[
    J^*
    =
    8x_1(18) + 4x_2(18)
    =
    8\left(
    \frac{50}{3}e^{36}
    -
    \frac{5}{3}e^{-18}
    \right)
    +
    4\left(
    \frac{50}{3}e^{36}
    +
    \frac{10}{3}e^{-18}
    \right)
    \]

    \[
    =
    \frac{400}{3}e^{36}
    -
    \frac{40}{3}e^{-18}
    +
    \frac{200}{3}e^{36}
    +
    \frac{40}{3}e^{-18}
    =
    \frac{600}{3}e^{36}
    =
    200e^{36}.
    \]   
\section*{Ejercicio 17.3}
Resolver
\[
\begin{aligned}
\min_{u} \quad & \int_0^1 u(t)\,dt \\
\text{subject to} \quad &
x'(t)=x(t)-u(t), \\
& x(0)=1,\quad x(1)=0, \\
& 1 \le u(t) \le 2.
\end{aligned}
\]

Tomando el hamiltoneano
\[
H(t,x,u)=  u+\la (x-u) = \la x +u(1-\la)
\]
donde su ecuación adjunta es 
 \[
 \la' = -\frac{\partial H}{\partial x} = -\la  \Rightarrow \la = ce^{-t}.
 \]
 Observe
 \[
 \psi (t) = \frac{\partial H}{\partial u} =1 -\la 
 \]
Dado que queremos minimizar el hamiltoneano vamos a analizar el signo de $\la$
\begin{itemize}
    \item si $1-\la <0$ el coeficiente en el hamiltoneano de u séra negativo  por lo tanto toca escoger el mayor $u*=2$
    \item  si $1-\la >0$ el coeficiente será positivo por tanto como buscamos minimizar escogemos $u*=1$
    \item si $1-\la = 0$ es un caso singular
\end{itemize}
Entonces tenemos que 
\begin{align*}
    u^*=1 \iff 1-\lambda>0 \iff 1-ce^{-t}>0 \iff t> \ln(c),\\
    u^*=2 \iff 1-\la <0 \iff  1-ce^{-t}<0 \iff t< \ln(c).
\end{align*}
por el resultado anterior tenemos que $c\leq0$ no puede darse por tanto vamos a ver que pasa si $0<c<1$ entonces
$1-ce^{-t}>0$ por lo tanto $u*=1$ asi $x'=x-1$ lo que nos da resultado  $x=ce{t}+1$ lo cual no cumple la condición de $x(0)=1$ por lo tanto
$c>1$
por lo tanto resolviendo por tramos para $0\leq t< \ln(c)$ tenemos que $x'=x-2$ donde nos queda 
\[
x=2-e^t
\] 
donde por la continuidad entonces en $t=\ln(c)$ tenemos que $x(\ln(c))=2-c$,
para el tramo de $\ln(c)<t\leq 1$ tenemos que $u=1$ y $x'=x-1$, lo que nos da $x=de^t+1$ y usando la condicion $x(\ln(c))=2-c$
tenemos que $d= \frac{1}{c}-1$
por lo tanto tenemos que 
\[
x(t)= (\frac{1}{c}-1)e^t+1
\]
y haciendo uso de $x(1)=0$ obtenemos que $c= \frac{e}{e-1}$, por lo tanto tenemos que 

\[
u^*(t)=
\begin{cases}
2, & 0 \le t \le t_s,\\[4pt]
1, & t_s < t \le 1,
\end{cases}
\qquad
\text{con } t_s = 1-\ln(e-1).
\]
con
    
\[
x^*(t)=
\begin{cases}
2-e^t, & 0 \le t \le t_s,\\[4pt]
1+\left(1-e^{t_s}\right)e^{\,t-t_s},
& t_s < t \le 1.
\end{cases}
\]
y el funcional nos da
\[
J_{\min}=2-\ln(e-1).
\]
\section*{Ejercicio 17.5}
Solucione 
\[
 \min \int_0^2 x(t)^2 dt
\]
Sujeto a $x'(t)=u(t), x(0)=0,x(2)=1, -1\leq u(t)\leq 1$.
Tenemos que el hamiltoneano generalizado $\lambda_0=1$ es 
\[
H(t,x,u,\lambda)= x^2+\la u
\]

\end{document}